\begin{afterword}
\summaryhead

The essay claims that three ideas are alike. A valid proof step never leads from a truth to a falsehood. A good reasoner judges an argument without regard to whether they like its conclusion. And law, which lets people with different interests escape bad equilibria, works only while people believe it is applied evenhandedly.

The author illustrates the second idea with two people who hear a hot day offered as evidence of global warming, and prefers the one who rejects it by symmetry with a cold day. He introduces the third with a juror in the Shkreli trial and with Democrats debating Al Franken's removal. He then argues that human law rests on rules that feel simple, that written law has decayed, and that impartiality is being lost to social media. A Chinese general and a judge in a 1962 novel serve as models. The essay ends by guessing that mathematicians are more honest than other people.

\responsehead

The essay has one sound idea, that an argument's validity is independent of whether its conclusion is true, and that idea is the textbook definition of validity. What the essay adds is the claim that this habit of mind, the fair enforcement of law, and the survival of civilization are one thing. It never argues that claim. The summary calls it a ``similarity,'' the title a ``key,'' and the last section admits that truth-seeking has ``no game theory and no morality.'' That admission leaves only a shared feeling of rule-following between the two halves. The essay spends about two-thirds of its length on law and never shows that the feeling does any work.

Judged by its own standard, the essay fails wherever the standard can be applied. Its model reasoner treats weak evidence as none, and by the essay's own Bayesian definition the reasoner it ranks lower is the more accurate one. It shortens the one opposing argument it links to until the argument is no longer the opponent's. The only person it shows giving up a bad argument for their own side is an imagined skeptic conceding a point to the author. Its closing hypothesis about mathematicians assumes the thesis it is meant to support.

The evidence is thin throughout, and where it can be checked it is often wrong. The Kelvin story is one that geophysicists have shown to be a myth, the Chinese general's son was a soldier from his home region, and the juror who opens the essay was removed by the procedure built to remove such jurors. The Alabama election offered as proof that the other side cooperates came after the Republican president and national committee had backed the accused candidate. Heinlein, the history of written law and the claim that no sane victim goes to the police come without sources. The one real case of unequal enforcement, racially skewed marijuana enforcement, gets a subordinate clause.

What the essay does for its reader is plainer than what it argues. It offers an identity, ``an intellectually strong mind,'' ``the unusually healthy of mind,'' and ties it to people who think like mathematicians. It tells the reader that the society around them is decaying and that they belong to the few who still keep The Law. The closing disclaimer of any ``grand agenda'' does not take back a title that promises the key to civilization.

In short: a definition from logic, joined to a theory of civilization by an analogy whose two halves, the essay concedes, work differently, supported by stories it did not check, and failing the standard it sets for everyone else.
\end{afterword}
